\chapter*{Resumen}
\addcontentsline{toc}{chapter}{Resumen}

Las eyecciones de masa coronal (CMEs) son los eventos solares más energéticos y de mayor influencia en el clima espacial, con una tasa de ocurrencia dependiente del ciclo solar. Cuando dos CMEs sucesivas ---una precursora y una sucesora--- interactúan en el medio interplanetario, la geoefectividad y el flujo de partículas energéticas solares se incrementan respecto a eventos individuales. Este trabajo simula numéricamente la interacción de CMEs homólogas, analizando su evolución cinemática y la deformación resultante de los perfiles de densidad y velocidad, con el fin de identificar umbrales de interacción. La metodología comprende: (i) el modelamiento de los perfiles de posición, velocidad y aceleración de cada CME mediante un modelo cinemático de dos fases; (ii) para una CME precursora con velocidad máxima $\sim 620$ km s$^{-1}$, la determinación de los umbrales de velocidad y aceleración que permiten la interacción con una sucesora, en un caso de aceleración abrupta (Caso 1) y uno de aceleración suave (Caso 2); y (iii) la evaluación bidimensional de los campos de densidad y velocidad resultantes para caracterizar el fortalecimiento del evento compuesto. Los resultados muestran que el umbral de interacción depende de la duración de la fase de aceleración de la sucesora, en el Caso 1 basta con $V_{\text{CME-2}}>0{,}86\,V_{\text{CME-1}}$, mientras que en el Caso 2 se requiere $V_{\text{CME-3}}>V_{\text{CME-1}}$. La CME resultante presenta densidad y velocidad superiores a las eyecciones individuales, y reproduce cualitativamente rasgos observacionales como la deformación morfológica progresiva de la densidad, lo que evidencia indirectamente un aumento en la geoefectividad y en el transporte de partículas energéticas, afectando el clima espacial.